\section{ماژول roast\_hashed\_linear (ROAST Hashed Linear)}
\label{sec:module_roast_hashed_linear}

\subsection{معرفی و انگیزه}
این مکانیزم بر مبنای مبانی نظری پیشرفته در ادبیات توجه و فشرده‌سازی \cite{ref_roasthashedlinear_2024} توسعه یافته است.

اشتراک‌گذاری پارامترها با استفاده از توابع درهم‌سازی تصادفی (\lr{Random Hashing}) رویکردی کارآمد برای کاهش ردپای حافظه ماتریس‌های وزن بدون نیاز به نگهداری ماتریس‌های کامل در فضای رم به شمار می‌رود. با این وجود، روش‌های سنتی مانند ترفند درهم‌سازی اسکالر (\lr{Scalar Hashing}) به دلیل عدم تطابق با الگوهای دسترسی حافظه در پردازنده‌های گرافیکی و نقض محلی بودن داده‌ها (\lr{Spatial Locality})، منجر به ناکارآمدی حافظه نهان و افت فاحش سرعت اجرای مدل می‌گردند.

ماژول \lr{roast\_hashed\_linear} با نام فارسی «لایه خطی درهم‌سازی‌شده با کاشی‌بندی ROAST» مبتنی بر ایده مقاله مرجع \lr{ROAST: A Hardware-Aware Weight Sharing Scheme for Efficient Neural Networks} توسعه یافته است. در این الگوریتم، وزن‌های شبکه به جای عناصر اسکالر منفرد، در سطح بلوک‌ها یا کاشی‌های دوبعدی (\lr{Tiling}) با جدول کوچکی از مقادیر به اشتراک گذاشته می‌شوند که سازگاری کاملی با سلسله‌مراتب حافظه سخت‌افزاری ایجاد می‌نماید.

\subsection{مبانی نظری و فرمول‌بندی ریاضی}
فرض کنید $\mathbf{W} \in \mathbb{R}^{d_{\text{out}} \times d_{\text{in}}}$ ماتریس وزن یک لایه خطی باشد. در الگوریتم \lr{ROAST}، ماتریس وزن به کاشی‌های مربعی با اندازه $B \times B$ افراز می‌شود:
\begin{equation}
T_{\text{out}} = \left\lceil \frac{d_{\text{out}}}{B} \right\rceil, \quad T_{\text{in}} = \left\lceil \frac{d_{\text{in}}}{B} \right\rceil
\end{equation}
تعداد کل کاشی‌ها برابر با $T_{\text{out}} \times T_{\text{in}}$ است. به جای ذخیره مقادیر مجزا، یک جدول وزن فشرده شامل $K$ باکت آموزش‌پذیر با پارامتر $\boldsymbol{\theta} \in \mathbb{R}^K$ تعبیه می‌شود ($K \ll d_{\text{out}} \cdot d_{\text{in}}$).

برای نگاشت هر کاشی در مختصات $(i, j)$ با $i \in \{0, \dots, T_{\text{out}}-1\}$ و $j \in \{0, \dots, T_{\text{in}}-1\}$ به یک باکت مشخص در جدول وزن، یک تابع درهم‌سازی با آنتروپی بالا و ترکیب بیتی (\lr{Bitwise Mixing}) فرمول‌بندی می‌شود:
\begin{align}
h_0 &= (i \cdot \mathtt{0x9E3779B1}) \oplus (j \cdot \mathtt{0x85EBCA77}) \\
h_1 &= h_0 \oplus (h_0 \gg 16) \\
h_2 &= h_1 \cdot \mathtt{0x7FEB352D} \\
h_3 &= (h_2 \oplus (h_2 \gg 15)) \ \& \ \mathtt{0x7FFFFFFF} \\
\operatorname{idx}(i, j) &= h_3 \bmod K
\end{align}
مقدار پارامتر کاشی $(i, j)$ از باکت استخراج می‌شود:
\begin{equation}
\mathbf{W}_{\text{tile}}(i, j) = \boldsymbol{\theta}[\operatorname{idx}(i, j)] \cdot \mathbf{1}_{B \times B}
\end{equation}
ماتریس کامل بازسازی‌شده با تکرار مقادیر کاشی‌ها تشکیل و به ابعاد مجاز برش داده می‌شود:
\begin{equation}
\hat{\mathbf{W}} = \left( \boldsymbol{\theta}[\operatorname{idx}(i, j)] \otimes \mathbf{1}_{B \times B} \right)_{:d_{\text{out}}, :d_{\text{in}}}
\end{equation}
همچنین در صورت استفاده از تطبیق کم‌مرتبه، اثر آداپتورهای \lr{LoRA} به صورت مستقیم در گذر پیشرو محاسبه می‌گردد:
\begin{equation}
\mathbf{y} = \mathbf{x} \hat{\mathbf{W}}^\top + \mathbf{b} + \frac{\alpha}{r} (\mathbf{x} \mathbf{A}_{\text{lora}}^\top) \mathbf{B}_{\text{lora}}^\top
\end{equation}

\paragraph{تحلیل طیفی و نامساوی‌های پایداری خطا}
برای بررسی میزان دقت تقریب ماتریس وزن یا نگاشت توجه در ماژول \lr{roast\_hashed\_linear}، تفاضل ماتریس فاکتورگیری‌شده از نگاشت متراکم اصلی بر اساس نامساوی هافدینگ-فون‌نیومان ارزیابی می‌گردد. فرض کنید $\mathbf{W} \in \mathbb{R}^{d_{\text{in}} \times d_{\text{out}}}$ ماتریس وزن هدف باشد که توسط تقریب کم‌مرتبه یا پراکنده با رتبه موثر $r$ جایگزین شده است. کران بالای خطای بازسازی بر مبنای نرم فروبنیوس به صورت زیر تصریح می‌شود:
\begin{equation}
\label{eq:spectral_bound_roast_hashed_linear}
\| \mathbf{W} - \mathbf{W}_{r} \|_F = \sqrt{\sum_{i=r+1}^{\min(d_{\text{in}}, d_{\text{out}})} \sigma_i^2(\mathbf{W})} \le \epsilon \| \mathbf{W} \|_F
\end{equation}
که در آن $\sigma_i(\mathbf{W})$ بیانگر $i$-امین مقدار منفرد مرتب‌شده ماتریس به صورت نزولی است. نسبت حفظ انرژی طیفی $\eta_r = \frac{\sum_{i=1}^r \sigma_i^2}{\sum_{i=1}^{d} \sigma_i^2}$ تضمین می‌کند که با انتخاب آستانه $\tau \ge 0.90$، بیش از نود درصد از پراکندگی واریانس داده‌ها حفظ شده و افت سیگنال در انتشار گرادیان رو به عقب به حداقل ممکن کاهش می‌یابد.

\paragraph{تحلیل مجانبی و پیچیدگی اشتراک حافظه سخت‌افزاری}
از منظر پیچیدگی محاسباتی و انتقال داده‌ها در سلسومراتب حافظه (حافظه سراسری \lr{HBM} به حافظه محلی \lr{SRAM})، عملیات ضرب ماتریسی در این مکانیزم نیازمند تخصیص بافر موقت با ابعاد فاکتورگیری‌شده است. برای پردازش دنباله‌ای به طول $N$ با بعد مخفی $D$، بار محاسباتی هسته از فرم درجه دو به رابطه زیر تقلیل می‌یابد:
\begin{equation}
\label{eq:complexity_roast_hashed_linear}
\mathcal{C}_{\text{roast\_hashed\_linear}} = \mathcal{O}\left( B \cdot N \cdot D \cdot r + B \cdot H \cdot N \cdot d_k \right) \ll \mathcal{O}\left( B \cdot H \cdot N^2 \cdot d_k \right)
\end{equation}
این بهبود مرتبه پیچیدگی، با کاهش چشمگیر حجم تبادل تانسورها در گذرگاه‌های ارتباطی کارت گرافیک، مانع از اشباع پهنای باند حافظه (پدیده \lr{Memory-Bound}) در فرآیند استنتاج طولانی‌مدت می‌گردد.

\subsection{معماری و پیاده‌سازی}
پیاده‌سازی این الگوریتم در فایل \lr{Codes/framework/compression\_engine/roast\_hashed\_linear.py} قرار گرفته و با شناسه‌های \lr{roast} و \lr{roast\_hashed\_linear} در رجیستری فشرده‌سازی سیستم موجود است. نقاط تمایز پیاده‌سازی عبارتند از:
\begin{itemize}
    \item \textbf{لایه برداری‌شده \lr{ROASTLinear}:} محاسبه ماتریس بازسازی‌شده با استفاده از عملگرهای سریع \lr{torch.repeat\_interleave} بدون هیچ‌گونه حلقه تکرار پایتونی پیاده‌سازی شده است.
    \item \textbf{حافظه نهان وزن بازسازی‌شده (\lr{Weight Caching}):} به منظور تسریع چشمگیر در زمان ارزیابی و استنتاج (\lr{eval mode})، وزن موثر بازسازی‌شده یک‌بار محاسبه و ذخیره موقت می‌شود تا از تکرار عملیات درهم‌سازی پرهیز گردد.
    \item \textbf{جدول بافت‌های تصادفی با قابلیت انجماد:} در زمان ریزتنظیم با \lr{LoRA}، جدول باکت‌های وزن منجمد گردیده و تنها ماتریس‌های کم‌مرتبه به‌روزرسانی می‌شوند تا همگرایی پایدار تضمین شود.
    \item \textbf{جایگزینی هوشمند با متد \lr{from\_linear}:} قابلیت دریافت وزن‌های لایه‌های چگال موجود، محاسبه میانگین بر روی کاشی‌ها، و مقداردهی متناسب جدول باکت‌ها به منظور تقریب اولیه مناسب وزن‌های پیش‌آموزش‌دیده.
\end{itemize}

\subsection{نتایج تجربی}
نتایج ارزیابی‌های تجربی ثبت‌شده برای ماژول \lr{roast\_hashed\_linear} در پایگاه اعتبارسنجی \lr{thesis\_final/master\_verified\_results.json} در جدول \ref{tab:benchmark_roast_hashed_linear} ارائه شده است:

\begin{table}[H]
\centering
\small
\caption{نتایج بنچ‌مارک تجربی فشرده‌ساز درهم‌ساز سخت‌افزار-آگاه \lr{roast\_hashed\_linear} روی مدل‌های پایه}
\label{tab:benchmark_roast_hashed_linear}
\begin{tabular}{lccccc}
\toprule
\textbf{مدل پایه \lr{(Base Model)}} & \textbf{دقت \lr{(Accuracy)}} & \textbf{شاخص \lr{F1}} & \textbf{حافظه \lr{(Peak VRAM)}} & \textbf{توان \lr{(Throughput)}} & \textbf{وضعیت} \\
\midrule
\lr{DistilBERT} & \lr{90.14\%} & \lr{0.9040} & \lr{337.59 MB} & \lr{54,047.60 tok/s} & \lr{PASS} \\
\lr{RoBERTa} & \lr{86.81\%} & \lr{0.8791} & \lr{583.49 MB} & \lr{28,100.97 tok/s} & \lr{PASS} \\
\lr{GPT-2} & \lr{51.72\%} & \lr{0.6774} & \lr{634.91 MB} & \lr{27,781.01 tok/s} & \lr{PASS\_LOW} \\
\lr{Qwen} & -- & -- & -- & -- & \lr{SKIP} \\
\bottomrule
\end{tabular}
\end{table}

بر اساس یافته‌های حاصل از بنچ‌مارک تجربی، ماژول \lr{roast\_hashed\_linear} پایداری و دقت بسیار بالایی در مدل‌های انکودر متراکم از خود نشان داده است؛ به طوری که در مدل \lr{DistilBERT} به دقت قابل توجه \lr{90.14\%} با F1 معادل \lr{0.9040} رسیده و در \lr{RoBERTa} دقت \lr{86.81\%} را احراز کرده است. در مدل خودبازگشت \lr{GPT-2} دقت برابر با \lr{51.72\%} و وضعیت \lr{PASS\_LOW} ثبت گردیده و در مدل \lr{Qwen} به علت ناسازگاری لایه‌های چندگروهی در نسخه ارزیابی اولیه، وضعیت به صورت \lr{SKIP} گزارش شده است.

\subsection{محدودیت‌ها و تحلیل شکست}
تحلیل نقاط ضعف و علل نوسان عملکرد در این ماژول به شرح زیر است:
\begin{itemize}
    \item \textbf{تصادم در درهم‌سازی کاشی‌ها (\lr{Hash Collisions}):} هنگامی که ابعاد ماتریس وزن بسیار بزرگ باشد، اشتراک‌گذاری ناگزیر چندین کاشی با رفتارهای آماری متفاوت در یک باکت واحد موجب نویز گرادیان و تداخل وزن‌ها می‌شود؛ این پدیده علت اصلی افت دقت در مدل \lr{GPT-2} ارزیابی شده است.
    \item \textbf{هزینه بازسازی برخط در غیاب کَشینگ سخت‌افزاری:} در محیط‌هایی با حافظه محدود که ذخیره ماتریس بازسازی‌شده امکان‌پذیر نیست، تکرار عملیات \lr{repeat\_interleave} در هر گام محاسباتی منجر به کاهش توان عملیاتی نسبت به لایه‌های استاندارد متراکم می‌گردد.
\end{itemize}
